% !TEX program = xelatex
\documentclass[12pt,a4paper]{article}

% Reusable Hebrew/RTL starting point for WORKFLOW-007.
% Replace placeholders and preserve user-specific macros when updating an existing artifact.

% ---------- Languages and fonts ----------
\usepackage{fontspec}
\usepackage{polyglossia}
\setdefaultlanguage{hebrew}
\setotherlanguage{english}

\setmainfont{Arimo}
\setsansfont{Arimo}
\newfontfamily\hebrewfont[Script=Hebrew]{Arimo}
\newfontfamily\hebrewfontsf[Script=Hebrew]{Arimo}
\newfontfamily\englishfont{Arimo}
\newfontfamily\englishfontsf{Arimo}

% ---------- Mathematics ----------
\usepackage{amsmath,amssymb,mathtools}

% ---------- Layout ----------
\usepackage[a4paper,top=2.0cm,bottom=2.0cm,left=2.2cm,right=2.2cm]{geometry}
\usepackage{setspace}
\setstretch{1.13}
\usepackage{microtype}
\usepackage{ragged2e}
\usepackage{needspace}
\usepackage{enumitem}
\usepackage{tabularx,booktabs,array}
\usepackage{xcolor}

\setlength{\parindent}{0pt}
\setlength{\parskip}{0.55em}
\setlength{\emergencystretch}{2em}
\setlist[itemize]{rightmargin=1.25em,leftmargin=0pt,itemsep=0.25em,topsep=0.3em}
\setlist[enumerate]{rightmargin=1.5em,leftmargin=0pt,itemsep=0.25em,topsep=0.3em}

% ---------- Colors ----------
\definecolor{navy}{HTML}{17324D}
\definecolor{bluegray}{HTML}{EAF1F7}
\definecolor{softred}{HTML}{FFF3F4}
\definecolor{deepred}{HTML}{8A2633}
\definecolor{softgreen}{HTML}{F2F8F4}
\definecolor{deepgreen}{HTML}{2D6245}
\definecolor{graytext}{HTML}{4E5963}
\definecolor{linegray}{HTML}{C9D2DA}

% ---------- Headings ----------
\newcommand{\runsection}[1]{%
  \clearpage
  \section*{\textcolor{navy}{#1}}%
  \addcontentsline{toc}{section}{#1}%
}

\newcommand{\runsubsection}[1]{%
  \Needspace{0.18\textheight}%
  \subsection*{\textcolor{navy}{#1}}%
  \addcontentsline{toc}{subsection}{#1}%
}

% ---------- Direction helpers ----------
\newcommand{\en}[1]{\textenglish{#1}}
\newcommand{\m}[1]{\LR{\ensuremath{#1}}}

% ---------- Reusable content units ----------
\newenvironment{contentunit}[1]{%
  \Needspace{0.27\textheight}%
  \par\medskip\noindent
  \begin{minipage}{\linewidth}
  \setlength{\parskip}{0.45em}%
  \colorbox{navy}{%
    \parbox{\dimexpr\linewidth-2\fboxsep\relax}{%
      \raggedleft\color{white}\bfseries #1%
    }%
  }%
  \par\vspace{0.45em}
}{%
  \par\vspace{0.2em}
  {\color{navy}\rule{\linewidth}{0.6pt}}%
  \end{minipage}
  \par\medskip
}

\newenvironment{rednote}[1]{%
  \Needspace{0.15\textheight}%
  \par\smallskip\noindent
  \begin{minipage}{\linewidth}
  \setlength{\parskip}{0.35em}%
  \colorbox{softred}{%
    \parbox{\dimexpr\linewidth-2\fboxsep\relax}{%
      \raggedleft\textcolor{deepred}{\bfseries #1}%
    }%
  }%
  \par\vspace{0.35em}
  \begingroup\leftskip=1em\rightskip=1em
}{%
  \par\endgroup
  {\color{deepred}\rule{\linewidth}{0.5pt}}%
  \end{minipage}
  \par\smallskip
}

\newenvironment{strengthnote}[1]{%
  \Needspace{0.15\textheight}%
  \par\smallskip\noindent
  \begin{minipage}{\linewidth}
  \setlength{\parskip}{0.35em}%
  \colorbox{softgreen}{%
    \parbox{\dimexpr\linewidth-2\fboxsep\relax}{%
      \raggedleft\textcolor{deepgreen}{\bfseries #1}%
    }%
  }%
  \par\vspace{0.35em}
  \begingroup\leftskip=1em\rightskip=1em
}{%
  \par\endgroup
  {\color{deepgreen}\rule{\linewidth}{0.5pt}}%
  \end{minipage}
  \par\smallskip
}

\newenvironment{noticebox}{%
  \Needspace{0.12\textheight}%
  \par\medskip\noindent
  \begin{minipage}{\linewidth}
  {\color{linegray}\rule{\linewidth}{0.6pt}}%
  \par\vspace{0.35em}
  \begingroup\leftskip=1em\rightskip=1em
}{%
  \par\endgroup\vspace{0.2em}
  {\color{linegray}\rule{\linewidth}{0.6pt}}%
  \end{minipage}
  \par\medskip
}

\newcommand{\keyrule}[1]{%
  \par\smallskip
  \begin{center}
    \fcolorbox{navy}{white}{%
      \parbox{0.9\linewidth}{\centering\bfseries\textcolor{navy}{#1}}%
    }
  \end{center}
  \smallskip
}

\pagestyle{plain}

\begin{document}
\begin{RTL}

% ============================================================
% Cover
% ============================================================
\thispagestyle{empty}
\vspace*{2.2cm}

\begin{center}
  {\fontsize{25}{30}\selectfont\bfseries\textcolor{navy}{דף שינונים ולקחים מצטבר}}\\[0.7em]
  {\Large \textbf{שם הקורס כאן}}\\[2.2em]

  \fcolorbox{navy}{bluegray}{%
    \parbox{0.72\linewidth}{%
      \centering
      \textbf{מבחנים שנכללו במסמך}\\[0.35em]
      שם מבחן ראשון; שם מבחן שני\\[0.75em]
      \textbf{גרסה:} 0.1 \qquad
      \textbf{עדכון אחרון:} תאריך
    }%
  }
\end{center}

\vfill
\begin{noticebox}
\textbf{מטרת המסמך:} לשמור רק לקחים בעלי ערך למבחן הבא: כללי זיהוי, טעויות שורש, בדיקות תקינות, שיטות קצרות, הצלחות שבריריות וחוזקות יציבות. זה אינו סיכום מלא של הקורס.
\end{noticebox}
\vfill

\begin{center}
\small\textcolor{graytext}{המסמך מיועד לקימפול באמצעות \en{XeLaTeX}.}
\end{center}

\clearpage
\tableofcontents
\clearpage

% ============================================================
% How to use
% ============================================================
\section*{\textcolor{navy}{איך משתמשים בדף}}
\addcontentsline{toc}{section}{איך משתמשים בדף}

\begin{contentunit}{סדר החזרה לפני מבחן}
\begin{enumerate}
  \item עוברים על מפת הזיהוי המהירה.
  \item קוראים יחידות שסומנו כפערים עמוקים או הצלחות שבריריות.
  \item עונים בעל פה על שאלות השינון.
  \item פותרים שאלת העברה אחת מכל משפחה חשובה.
\end{enumerate}
\end{contentunit}

% ============================================================
% Recognition map
% ============================================================
\runsection{מפת זיהוי מהירה}

\begingroup
\hbadness=10000
\renewcommand{\arraystretch}{1.35}
\small
\begin{center}
\begin{tabularx}{0.97\textwidth}{>{\raggedleft\arraybackslash}p{0.25\textwidth} >{\raggedleft\arraybackslash}p{0.30\textwidth} >{\raggedleft\arraybackslash}X}
\toprule
\textbf{מה רואים בשאלה} & \textbf{הכלי הראשון} & \textbf{הבדיקה שאסור לדלג עליה} \\
\midrule
תיאור משפחת שאלה & הכלי או הייצוג הראשון & תנאי, תחום, סימן או בדיקת סבירות \\
תיאור משפחת שאלה נוספת & הכלי או הייצוג הראשון & הבדיקה המכריעה \\
\bottomrule
\end{tabularx}
\end{center}
\normalsize
\endgroup

% ============================================================
% Exam-specific section template
% ============================================================
\runsection{מבחן 1 — שם המבחן}

\begin{noticebox}
\textbf{מעמד התוכן:} ציינו אילו מסקנות נבדקו מול פתרון רשמי, אילו שוחזרו ממקורות הקורס ואילו הן הסבר עצמאי.
\end{noticebox}

\begin{contentunit}{שם המושג או משפחת השאלה}
\textbf{הפער שהתגלָה:} תיאור מדויק וקצר של נקודת הסטייה הראשונה או של ההצלחה השברירית.

\textbf{כלל או משפט:} כתבו את הכלל המדויק ואת תנאי השימוש בו.

\textbf{סדר הפעולות:}
\begin{enumerate}
  \item שלב זיהוי.
  \item שלב בחירת השיטה.
  \item שלב חישוב או הוכחה.
  \item בדיקת תקינות.
\end{enumerate}

\keyrule{כלל קצר, כללי וזכיר שאינו תלוי במספרים של השאלה המקורית.}
\end{contentunit}

\begin{rednote}{מלכודת או בדיקת תקינות}
כתבו מה היה יכול לעצור את הטעות בזמן: טווח הסתברות, סימן, תחום תמיכה, יחידות, תנאי משפט או מקרה קצה.
\end{rednote}

\begin{strengthnote}{חוזקה יציבה או הצלחה שברירית}
כתבו את הראיה החיובית. אם התשובה הייתה נכונה אך הדרך לא יציבה, ציינו זאת במפורש וכתבו מה עדיין צריך לחזק.
\end{strengthnote}

% ============================================================
% Cumulative sections
% ============================================================
\runsection{לקחים מצטברים ובדיקות תקינות}

\begin{contentunit}{שאלת הפתיחה לפני חישוב}
\begin{enumerate}
  \item מה המשתנה או הביטוי סופר, מודד או מתאר?
  \item איזה ייצוג הופך את האירוע או הכמות לפשוטים?
  \item מה תחום התמיכה או תנאי התקפות?
  \item איזו בדיקת סבירות יכולה לפסול תוצאה בלתי אפשרית?
\end{enumerate}
\end{contentunit}

\runsection{ידע שכבר עובד וצריך לשמר}

\begin{strengthnote}{שם החוזקה}
תארו ראיות חוזרות או ראיה חזקה לכך שהזיהוי, השיטה והביצוע יציבים.
\end{strengthnote}

\runsection{שאלות שינון מהירות}

\begin{contentunit}{שם הנושא}
\textbf{שאלה:} שאלה קצרה שבודקת את כלל הזיהוי, התנאי או בדיקת התקינות.\\
\textbf{תשובה:} תשובה קצרה ומדויקת.
\end{contentunit}

% ============================================================
% Future update protocol
% ============================================================
\runsection{תבנית לעדכון לאחר מבחן נוסף}

\begin{contentunit}{מבחן חדש}
לכל מבחן עתידי מוסיפים:
\begin{itemize}
  \item שם המבחן ותאריך הבדיקה;
  \item מקור רשמי או מעמד הפתרון;
  \item טעויות שורש והשלכות נגררות;
  \item אי-הבנות שהתגלו בשיחה שלאחר הבדיקה;
  \item הצלחות שבריריות וחוזקות יציבות;
  \item בדיקות תקינות שהיה כדאי להפעיל;
  \item שאלות שינון חדשות;
  \item עדכון מפת הזיהוי רק כאשר נוסף כלל כללי בעל ערך חוזר.
\end{itemize}
\end{contentunit}

\end{RTL}
\end{document}
